% from ega4-2-fr.tex
\oldpage[IV-2]{118}

Posons pour abréger $\mathcal F'=\mathcal H^0_{X/Z}(\mathcal F)$. On sait
(5.9.11) que $\mathcal F'$ est $Z$-clos et l'assertion (i) résulte donc de
(5.10.5) appliqué à $\mathcal F'$. Soit
$x\in\operatorname{Ass}(\mathcal F)\cap(X-Z)$; comme les restrictions de
$\mathcal F$ et $\mathcal F'$ à $X-Z$ sont canoniquement isomorphes (5.9.11),
on a $x\in\operatorname{Ass}(\mathcal F')$. Considérons un point
$y\in Z\cap\overline{\{x\}}$; l'idéal premier $\mathfrak p$ de
$\mathcal O_y$ correspondant à $x$ est associé au $\mathcal O_y$-module
$\mathcal F'_y$, donc on a, en vertu de (i) et de (0, 16.4.6.2),
\[
2\leq\operatorname{prof}(\mathcal F'_y)
\leq\dim(\mathcal O_y/\mathfrak p)
=\operatorname{codim}(\overline{\{y\}},\overline{\{x\}}),
\]
d'où (ii). Enfin, pour prouver (iii), notons que $U$ est l'ensemble des
$x\in X$ tels que $\widetilde{\mathcal F}_x$ soit $Z_x$-clos (5.10.5), ou
encore, en vertu de (5.9.6), l'ensemble des points où l'homomorphisme
canonique $\mathcal F_x\to\mathcal F'_x$ est bijectif; c'est donc le
complémentaire de la réunion des supports de $\operatorname{Ker}(\rho_{X/Z})$
et $\operatorname{Coker}(\rho_{X/Z})$, et ces derniers sont des
$\mathcal O_X$-Modules cohérents en vertu de l'hypothèse (0$_\mathrm I$,
5.3.4) donc ont un support fermé (0$_\mathrm I$, 5.2.2); cela montre que $U$
est ouvert, et $U$ est évidemment le plus grand ouvert tel que $\mathcal F|U$
soit $(Z\cap U)$-clos; enfin l'inclusion $X-U\subset Z$ résulte de (5.9.11).

Nous verrons plus loin (5.11.1) que dans les cas les plus importants
l'assertion (ii) entraîne réciproquement que
$\mathcal H^0_{X/Z}(\mathcal F)$ est cohérent.

\begin{env}[5.10.11]
\label{IV.2.s5.10.11-fr}
Soient $X$ un préschéma localement noethérien, $Z$ une partie de $X$ stable
par spécialisation; on a vu que
$\mathcal A=\mathcal H^0_{X/Z}(\mathcal O_X)$ est une
$\mathcal O_X$-Algèbre quasi-cohérente (5.9.3); le $X$-schéma
$X'=\operatorname{Spec}(\mathcal H^0_{X/Z}(\mathcal O_X))$ (II, 1.3.1) est
appelé la $Z$-clôture de $X$. En outre pour tout $\mathcal O_X$-Module
$\mathcal F$, $\mathcal H^0_{X/Z}(\mathcal F)$ est un $\mathcal A$-Module,
qui est quasi-cohérent si $\mathcal F$ est un $\mathcal O_X$-Module
quasi-cohérent; dans ce dernier cas, il y a donc un unique
$\mathcal O_{X'}$-Module $\mathcal F'$ tel que l'on ait, en désignant par
$g:X'\to X$ le morphisme structural (II, 1.4.3)
\begin{equation}
\mathcal H^0_{X/Z}(\mathcal F)=g_*(\mathcal F').
\tag{5.10.11.1}\label{IV.2.s5.10.11.1-fr}
\end{equation}
\end{env}

\begin{proposition}[5.10.12]
\label{IV.2.s5.10.12-fr}
Les notations étant celles de (5.10.11):
\begin{enumerate}
\item[(i)] Soit $x$ un point de $X$. Pour que le morphisme
\[
X'\times_X X_x\longrightarrow X_x(=\operatorname{Spec}(\mathcal O_{X,x}))
\]
déduit de $g$ par localisation soit un isomorphisme, il faut et il suffit que
$\mathcal O_X$ soit $Z$-clos au point $x$ (ce qui a lieu pour tout $x\in X-Z$).
\item[(ii)] Posons $Z'=g^{-1}(Z)$, et supposons $X'$ localement noethérien.
Alors $\mathcal F'$ est $Z'$-clos; si de plus $\mathcal F'$ est un
$\mathcal O_{X'}$-Module cohérent, on a
$\operatorname{prof}_{Z'}(\mathcal F')\geq2$.
\item[(iii)] Supposons que $\mathcal H^0_{X/Z}(\mathcal O_X)$ et
$\mathcal H^0_{X/Z}(\mathcal F)$ soient cohérents. Alors le morphisme
$g:X'\to X$ est fini; l'ensemble $U$ des $x\in X$ tels que l'on ait
$\operatorname{prof}_{Z_x}(\widetilde{\mathcal O}_x)\geq2$ et
$\operatorname{prof}_{Z_x}(\widetilde{\mathcal F}_x)\geq2$ est ouvert dans
$X$ et tel que $X-U\subset Z$; en outre $U$ est le plus grand ensemble ouvert
de $X$ tel que la restriction $g^{-1}(U)\to U$ de $g$ soit un isomorphisme et
que la restriction $\mathcal F|U\to\mathcal F'|g^{-1}(U)$ du $g$-morphisme
canonique $\mathcal F\to\mathcal F'$ soit un isomorphisme.
\end{enumerate}
\end{proposition}

L'assertion (i) résulte des définitions, et (iii) est une conséquence immédiate
de (5.10.10, (iii)). Pour prouver (ii), considérons un ouvert $V$ de $X$
contenant $Z-X$ et son image réciproque $V'=g^{-1}(V)$; si $i:V\to X$ et
$i':V'\to X'$ sont les injections

